
% SGA 3，第 VII 篇_B；英文重建本第一轮的中文译本。
% 范围： 节 2.9 全文, 包括命题 2.9.1 与
% 注 2.9.2, 依据文献局部页 65--66.
% 止于下节之前： 节 3.
% 合并阅读本页： 578--579. 重版印刷页： 567--568.
% 范围内编者注： 118.
% 依据： Polo--Gille Exp7B-13oct24.pdf,
% SHA-256 D303690E0195654BBE7C3DAA5D89EA23E24CD709A3575217234AC19AE8836694.
% 参照了 Jacob Reinhold 的 CC BY 4.0 Markdown，作为对照来源。
% 本组件无图；所有数学内容均为原生 LaTeX。
\subsection*{2.9. 域上的点化余交换 Hopf 代数}
\phantomsection
\label{sga3:VIIB:2.9}
\addcontentsline{toc}{subsection}{\SGARefLink{sga3:VIIB:2.9}{2.9}. 域上的点化余交换 Hopf 代数}
\footnote{\label{sga3:VIIB:nde:118}编者添加此小节，以余交换 Hopf 代数的语言表述命题 2.5.2 及编者注 100 中指出的变体。}
设 \(k\) 是域，\(C\) 是 \(k\)-余代数，并赋予
\[
  A=C^*
\]
以 \SGARefLink{sga3:VIIB:2.8}{2.8} 所描述的不一定交换的预有限 \(k\)-代数结构。由 \SGARefLink{sga3:VIIB:0.2.2}{0.2.2}，
\[
  C=A^\dagger=\operatorname{Hom}_c(A,k).
\]
因此映射
\[
  D\longmapsto
  D^\perp=\{f\in A\mid f(D)=0\}
\]
是从 \(C\) 的子余代数到 \(A\) 的闭双边理想\textup{（以下将简称为理想）}的双射。其逆为
\[
  I\longmapsto
  I^\perp
  =
  \{x\in C=A^\dagger\mid x(I)=0\}.
\]
由于每个极大闭理想都是极大开理想\textup{（参见\ \SGARefLink{sga3:VIIB:0.2.1}{0.2.1}）}，每个子余代数都包含一个单子余代数，后者必为有限维。
回顾，若子余代数 \(D\subset C\) 包含唯一的单子余代数 \(S_0\)，则称 \(D\) 是\emph{不可约的}。等价地，
\[
  \mathfrak m_0=S_0^\perp
\]
是包含 \(D^\perp\) 的唯一开极大理想；又等价地，
\[
  D^\perp+\mathfrak m=A
  \qquad
  \text{对每个开极大理想 }
  \mathfrak m\ne\mathfrak m_0.
\]
特别地，这适用于 \(D=S_0\)。
令 \(\Sigma_0\) 为包含 \(S_0\) 的所有不可约子余代数 \(C_i\) 的和。显然它是子余代数，且是不可约的，因为对每个 \(\mathfrak m\ne\mathfrak m_0\)，\SGARefLink{sga3:VIIB:0.2.D}{0.2.D} 给出
\[
  \mathfrak m+\bigcap_i C_i^\perp
  =
  \bigcap_i\bigl(\mathfrak m+C_i^\perp\bigr)
  =
  A.
\]
称余代数 \(\Sigma_0\) 为 \emph{\(C\) 中对应于 \(S_0\) 的不可约分量}。
\textit{法文在最后一句印作 \(C_0\)；根据前面的定义，所指的单子余代数应为 \(S_0\)。}
若余代数 \(C\) 的每个单 \(k\)-子余代数都是一维的，则称 \(C\) 是\emph{点化的}。等价地，对 \(A\) 的每个开极大理想 \(\mathfrak m\)，
\[
  A/\mathfrak m=k
\]
成立。还要回顾，若 \(C\) 是不可约且点化的，则称其为\emph{连通的}。若 \(C\) 为双代数，则它连通当且仅当它不可约，因为 \(k\cdot1_C\) 是单子余代数。
此后设 \(C\) 是余交换的。则 \(A\) 是交换的，因而是其局部分量的积
\[
  A_{\mathfrak m},
  \qquad
  \mathfrak m\in\Upsilon(A);
\]
参见\ \SGARefLink{sga3:VIIB:0.1.1}{0.1.1}。置
\[
  S_{\mathfrak m}
  =
  \mathfrak m^\perp
  \simeq
  (A/\mathfrak m)^*,
\]
并令 \(\Sigma_{\mathfrak m}\) 为其不可约分量。可以明确描述这个分量。令
\[
  J_{\mathfrak m}
  =
  \ker\bigl(A\longrightarrow A_{\mathfrak m}\bigr).
\]
则 \(J_{\mathfrak m}\subset\mathfrak m\)，而且它是满足
\[
  I+\mathfrak m'=A
  \qquad
  \text{对每个开极大理想 }
  \mathfrak m'\ne\mathfrak m.
\]
的最小闭理想 \(I\subset\mathfrak m\)。事实上，每个这样的 \(I\) 都包含所有 \(\mathfrak m'\ne\mathfrak m\) 对应的 \(A_{\mathfrak m'}\)，因而包含 \(J_{\mathfrak m}\)。由于
\[
  A=J_{\mathfrak m}\oplus A_{\mathfrak m},
\]
得到
\[
  \Sigma_{\mathfrak m}
  =
  J_{\mathfrak m}^{\perp}
  \simeq
  A_{\mathfrak m}^{\dagger}.
\]
现在可以证明：
\medskip
\noindent\textbf{命题 2.9.1。}
\phantomsection
\label{sga3:VIIB:2.9.1}
{\itshape 设 \(k\) 是域。
\begin{enumerate}
\item[\textup{\SGAInlineRefTarget{sga3:VIIB:2.9.1:item:i}(i)}]
设 \(G\) 是所有剩余域均等于 \(k\) 的形式 \(k\)-群。则 \(G_e\) 是常值 \(k\)-群 \(M_k\)，其中
\[
  M
  =
  G(k)
  =
  \left\{
    x\in\mathbf H(G)
    \,\middle|\,
    \varepsilon(x)=1,\quad
    \Delta(x)=x\otimes x
  \right\},
\]
且 \(\mathbf H(G)\) 是 \(\mathbf H(G_0)\) 与群代数 \(kM\) 的半直积\textup{（参见\ \SGARefLink{sga3:VIIB:2.4.B}{2.4.B}）}。
\item[\textup{\SGAInlineRefTarget{sga3:VIIB:2.9.1:item:ii}(ii)}]
等价地，设 \(H\) 是点化余交换 Hopf \(k\)-代数。则 \(H\) 是单位元 \(1_H\) 的不可约分量 \(H_0\) 与 \(kM\) 的半直积，其中
\[
  M
  =
  \{x\in H\mid \varepsilon(x)=1,\ \Delta(x)=x\otimes x\}.
\]
\end{enumerate}}
\textit{证明。}
先证明\textup{\SGARefLink{sga3:VIIB:2.9.1:item:i}{(i)}}。由于 \(G\) 的每个剩余域都等于 \(k\)，投影
\[
  \pi:G\longrightarrow G_e
\]
有截面 \(s:G_e\to G\)，在每个点 \(g\in G\) 处定义为
\[
  \mathcal O_{G,g}\longrightarrow\kappa(g)=k.
\]
此外，对所有 \(g,h\in G\)，代数
\[
  \mathcal O_{G,g}
  \widehat\otimes_k
  \mathcal O_{G,h}
\]
是剩余域为 \(k\) 的局部代数。因此 \(s\times s\) 是
\[
  \pi\times\pi:
  G\times G\longrightarrow
  (G\times G)_e=G_e\times G_e.
\]
的截面。由于 \(\pi\) 是群态射，
\[
  \pi\circ m_G\circ(s\times s)
  =
  m_{G_e}
  =
  \pi\circ s\circ m_{G_e}.
\]
由于 \(\pi\) 是满态射，可知 \(s\) 是群态射。故
\[
  G=G_0\rtimes G_e,
\]
且 \(\mathbf H(G)\) 是 \(\mathbf H(G_0)\) 与 \(\mathbf H(G_e)\) 的半直积。
由于 \(G\) 的每个剩余域都为 \(k\)，\(\mathbf H(G_e)\) 是群代数 \(kM\)，其中
\[
  M=G_e(k);
\]
参见\ \SGARefLink{sga3:VIIB:2.5.A}{2.5.A}。最后，由于 \(G_0\) 无穷小，映射
\[
  G(k)\longrightarrow G_e(k)
\]
是单射。它也是满射，因为它有由 \(s\) 诱导的截面。因此 \(M=G(k)\)，证明了\textup{\SGARefLink{sga3:VIIB:2.9.1:item:i}{(i)}}。
对于\textup{\SGARefLink{sga3:VIIB:2.9.1:item:ii}{(ii)}}，只须再证明
\[
  H_0=\mathbf H(G_0).
\]
\(H\) 的单位元 \(1_H\) 是增广
\[
  \varepsilon_{\mathcal A}:\mathcal A\longrightarrow k,
\]
它是 \(G(k)\) 的单位截面。因此，仿射代数 \(\mathcal A(G)\) 中对应于 \(H_0\) 的局部分量恰好是 \(\mathcal A(G_0)\)。由上述描述，
\[
  H_0
  =
  \mathcal A(G_0)^\dagger
  =
  \mathbf H(G_0).
\]
命题得证。
\medskip
\noindent\textbf{注 2.9.2。}
\phantomsection
\label{sga3:VIIB:2.9.2}
\begin{enumerate}
\item[\textup{(a)}]
上述命题隐含于 \SGARefLink{sga3:VIIB:2.5.2:proposition}{2.5.2}，B.\ Kostant 也曾独立得到它\textup{（参见 [Sw69]，序言）}。将其与下文 Cartier 的\SGARefLink{sga3:VIIB:3.3:cartier-theorem}{定理 3.3}\textup{（参见 [Ca62]，\S~12，定理 3）}结合——该定理 Kostant 也曾得到\textup{（参见 [Sw69]，同上）}——所得结果常称为 \emph{Cartier--Gabriel--Kostant 定理}。
\item[\textup{(b)}]
R.\ G.\ Heyneman 与 M.\ E.\ Sweedler 将\SGARefLink{sga3:VIIB:2.9.1}{命题 2.9.1} 的形式\textup{\SGARefLink{sga3:VIIB:2.9.1:item:ii}{(ii)}} 推广到假设 \(H\) 点化且是其不可约分量的直和、但不一定余交换的情形；参见 [HS69]，定理 3.5.8。
\end{enumerate}
